% =============================================================================
% Version 3 (simplified rewrite).
%
% Build:  pdflatex paper && bibtex paper && pdflatex paper && pdflatex paper
%
% Figures resolve through the "images" symlink -> ../shared/figures
% Bibliography resolves through the "references.bib" symlink -> ../shared/references.bib
% =============================================================================

\documentclass[10pt,a4paper]{article}

% ---------------------------------------------------------------- packages ---
\usepackage[top=2.2cm,bottom=2.2cm,left=2.3cm,right=2.3cm]{geometry}
\usepackage{float}
\usepackage{subcaption}
\usepackage{booktabs}
\usepackage[T1]{fontenc}
\usepackage{newtxtext,newtxmath}
\usepackage{graphicx}
\usepackage[numbers,sort&compress]{natbib}
\usepackage[hidelinks]{hyperref}

% ---------------------------------------------------------------- layout -----
\setlength{\parindent}{0pt}
\setlength{\parskip}{0.4em}
\pagestyle{plain}

\begin{document}

% ---------------------------------------------------------------- title ------
{\LARGE\bfseries
Explainable Survival Prediction from Gene Expression and Clinical Data
\par}

\vspace{1em}

{\large
xyz$^{a}$,
xyz$^{b}$,
xyz$^{a,*}$,
xyz$^{a}$,
xyz$^{a}$,
xyz$^{b}$,
xyz$^{a}$
\par}

\vspace{0.8em}

{\small
$^{a}$ Department of Computer Science and Engineering,
Amity University Kolkata, Kolkata, West Bengal, India

$^{b}$ Department of Computational Biology,
Indraprastha Institute of Information Technology Delhi,
New Delhi, India

$^{*}$ Corresponding author.
}

% ---------------------------------------------------------------- abstract ---
\vspace{1.2em}

{\bfseries ABSTRACT\par}

Doctors usually estimate how long a cancer patient will live from the size and
spread of the tumour. Two patients placed in the same group often live for very
different lengths of time. Measurements of gene activity can explain part of
that difference. Models built on such measurements are usually hard to read, so
a model was built here that stays readable at every step. Clinical records and
gene activity were combined for 1,620 patients from four public cancer datasets.
The datasets cover brain, liver, pancreatic and skin cancer. Highly similar
genes were removed first, which reduced 500 candidate genes to 59. These 59
genes were compressed into 50 summary variables using principal component
analysis, a standard method for reducing the number of variables. A Cox
proportional hazards model was then fitted with an elastic-net penalty, which
shrinks unhelpful coefficients towards zero. On a held-out test set of 324
patients the model reached a concordance index of 0.736. The concordance index
measures how often two patients are ranked in the correct survival order. Raw
risk scores were converted into survival probabilities using isotonic
regression. Brier scores, which measure probability error, ranged from 0.132 to
0.178 across the 12, 24, 36 and 60-month horizons. Uncertainty for each patient
was estimated from 5,000 simulated survival times. The compression step was
reversed by matrix multiplication, which returned an exact risk weight for every
one of the 59 genes. No approximate explanation method was needed.

% ---------------------------------------------------------------- keywords ---
\vspace{1em}
{\bfseries KEYWORDS\par}

Survival analysis; gene expression; Cox proportional hazards; elastic net;
probability calibration; model explainability

\vspace{3pt}

% =============================================================================
\section{Introduction}
% =============================================================================

Cancer treatment decisions depend on an estimate of how long a patient is likely
to live. That estimate is usually built from the tumour, node and metastasis
(TNM) staging system \citep{amin2017ajcc}. The system groups patients by tumour
size and by how far the disease has spread. Staging describes anatomy. It does not
describe the molecular activity inside the tumour. Two patients placed in the
same stage therefore often survive for very different lengths of time. Part of
that gap can be traced to differences in gene activity. Ribonucleic acid
sequencing (RNA-Seq) measures how active each gene is in a tissue sample
\citep{li2011rsem}. Some genes measured this way are already linked to poor
outcomes. High activity of the gene CXCL5 is associated with shorter survival in
liver cancer \citep{zhou2012overexpression}. High activity of the gene CA9
marks tumour regions starved of oxygen, and is also linked to worse outcomes
\citep{giatromanolaki2001carbonic}. The Cancer Genome Atlas (TCGA) has released
matched clinical records and gene activity measurements for dozens of cancer
types \citep{weinstein2013cancer}. This raises a direct question. Can staging
information and gene activity be combined into one model that predicts survival
better than staging alone, and that a doctor can still inspect?

The standard tool for this task is the Cox proportional hazards model
\citep{cox1972regression}. The model works with a hazard, which is the
instantaneous risk of dying at a given moment for a patient who is still alive.
The Cox model does not assume any particular shape for how that baseline risk
changes over time. It only assumes that each patient's hazard is the baseline
hazard multiplied by a fixed factor. That factor comes from a weighted sum of
the patient's measurements. The model also handles right censoring, which occurs
when a patient is still alive when the study ends. Such a patient contributes
the information that they survived at least that long. This flexibility comes
with a limit. The number of measurements must stay well below the number of
patients. When the number of measurements approaches the number of patients, the
fitted weights are no longer uniquely determined \citep{tibshirani1997lasso}.
Gene activity data breaks this limit immediately, because a single experiment
measures thousands of genes for a few hundred patients.

Three families of methods have been used to work around this limit. The first
family adds a penalty term to the Cox model that pushes weights towards zero.
The lasso penalty sets some weights exactly to zero and so selects a subset of
genes \citep{tibshirani1997lasso}. The elastic-net penalty adds a second
term to that behaviour \citep{zou2005regularization, wu2012elastic}. The second
term shrinks correlated genes together, instead of picking one gene from each
correlated group at random \citep{simon2011regularization}. Standard textbook
treatments of both penalties are available \citep{hastie2009elements}. The
second family reduces the number of variables before fitting. Principal
component analysis (PCA) replaces many correlated measurements with a few
summary variables that capture most of the variation
\citep{hotelling1933analysis, jolliffe2016principal}. Cross-validated penalised
Cox models applied to gene activity data follow this pattern
\citep{vanhouwelingen2006crossvalidated}. The third family corrects the
probabilities a model reports. Isotonic regression maps a raw score onto a
probability \citep{ayer1955empirical}. It preserves the ordering of patients and
assumes nothing about the shape of that mapping \citep{zadrozny2002transforming}.
Separately, neural network models such as DeepSurv \citep{katzman2018deepsurv},
Cox-nnet \citep{ching2018coxnnet} and Cox-PASNet \citep{hao2019coxpasnet}
replace the weighted sum with a flexible function. These models can capture
effects that a weighted sum cannot. They also need more patients than a single
cancer cohort usually provides, and their internal weights no longer map onto
named genes.

Four specific problems remain unsolved when these methods are used together.
First, groups of genes that act in the same biological pathway are often left in
the model together. Their measurements are almost identical, which makes the
fitted weights unstable across repeated fits \citep{simon2011regularization}.
Second, the Cox model reports a relative risk score. A doctor cannot read a
one-year or five-year survival probability off that score without a separate
conversion step, and that step is frequently skipped. Third, a single predicted
number says nothing about the range of outcomes that are plausible for one
patient. That range matters for surgical and palliative care planning. Fourth,
PCA solves the variable-count problem but creates a readability problem. The
fitted weights attach to summary variables, not to named genes. Post-hoc
explanation methods such as SHapley Additive exPlanations (SHAP)
\citep{lundberg2017unified} and Local Interpretable Model-agnostic Explanations
(LIME) \citep{ribeiro2016why} estimate gene-level importance by sampling. They
approximate the answer rather than recovering it exactly.

This work addresses all four problems in a single model. Five contributions are
reported. A correlation filter removes redundant genes using the training data
only, reducing 500 candidate genes to 59. A Cox model with an elastic-net
penalty is fitted using a smoothed objective, which keeps the optimiser stable
where the penalty is otherwise not differentiable. Isotonic regression
calibrators, fitted on a separate split, convert risk scores into survival
probabilities at four time horizons. A simulation draws 5,000 possible survival
times per patient, giving pessimistic, median and optimistic estimates. Finally,
the PCA step is reversed by one matrix multiplication. This returns an exact risk
weight for each of the 59 genes. It also returns an exact additive breakdown for
each patient. Figure~\ref{fig:unified_architecture} shows how data moves through
these stages.

\begin{figure}[H]
    \centering
    \includegraphics[width=0.85\textwidth]{images/fig1_unified_architecture.png}
    \caption{Overall data flow. Raw clinical records and gene activity
    measurements enter on the left. They pass through the correlation filter,
    the compression step, the penalised Cox model, the probability calibrator
    and the simulation stage. Gene-level explanations leave on the right.}
    \label{fig:unified_architecture}
\end{figure}

% =============================================================================
\section{Literature Survey}
% =============================================================================

\subsection{Statistical models}

Early work on gene activity and survival extended the Cox model with penalty
terms. Ridge, lasso and elastic-net penalties were all applied to the Cox
partial likelihood \citep{tibshirani1997lasso, zou2005regularization,
simon2011regularization, wu2012elastic}. These methods keep the weighted-sum
structure of the Cox model, so each fitted weight still refers to one named
measurement. Cross-validated penalised Cox regression on microarray data
established the practical recipe still in use
\citep{vanhouwelingen2006crossvalidated}. That recipe fits the penalty strength
by cross-validation. Performance is then reported on data the model never saw.
Random survival forests took a different route \citep{ishwaran2008random}. They
drop the assumption that each patient's hazard is a fixed multiple of the
baseline hazard. Error is estimated from the samples left out of each tree.

\subsection{Neural network models}

DeepSurv replaced the weighted sum in the Cox model with a neural network
\citep{katzman2018deepsurv}. Cox-nnet applied the same idea to gene activity
data from ten TCGA cohorts \citep{ching2018coxnnet}. Cox-PASNet added biological
pathway structure to the network, and was tested on TCGA brain and ovarian
cancer data \citep{hao2019coxpasnet}. DeepProg combined several deep and
classical models into an ensemble across many TCGA cancer types
\citep{poirion2021deepprog}. It is notable for testing on independent cohorts
rather than by cross-validation alone. Dynamic-DeepHit went further
\citep{lee2020dynamic}. It uses repeated measurements over time to update its
predictions as new data arrives. Across this group, the concordance index is almost
always the headline metric. Calibrated probabilities, per-patient uncertainty
and external validation are reported far less consistently.

\subsection{Feature representation}

Most models in this area use gene activity data together with clinical
variables. Some add measurements of microRNA, deoxyribonucleic acid methylation
or tissue images. Autoencoder models compress several such data types into a
small set of learned variables. Chaudhary and colleagues applied this approach
to 360 TCGA liver cancer patients \citep{chaudhary2018deep}. They reported a
concordance index of 0.68, and tested on independent liver cancer cohorts. Learned
variables of this kind can capture patterns a weighted sum cannot. They are also
difficult to trace back to individual genes, which is the readability gap this
work targets.

\subsection{Comparison of existing models}

Table~\ref{tab:comparative_models} summarises the studies discussed above.

\begin{table}[htbp]
\centering
\caption{Summary of related survival prediction studies.}
\label{tab:comparative_models}
\scriptsize
\setlength{\tabcolsep}{2.5pt}
\renewcommand{\arraystretch}{1.15}
\resizebox{\columnwidth}{!}{%
\begin{tabular}{p{2.2cm} p{2.4cm} p{2.4cm} p{2.2cm} p{2.6cm}}
\toprule
\textbf{Study} &
\textbf{Cohort} &
\textbf{Input data} &
\textbf{Model} &
\textbf{Validation and metric} \\
\midrule

van Houwelingen et al. \citep{vanhouwelingen2006crossvalidated} &
Breast cancer, microarray &
Gene expression &
Penalised Cox &
Cross-validation \\

Ishwaran et al. \citep{ishwaran2008random} &
General survival data &
Clinical and mixed &
Random survival forest &
Out-of-bag error \\

Katzman et al. \citep{katzman2018deepsurv} &
Clinical cohorts &
Clinical data &
DeepSurv &
Concordance index \\

Ching et al. \citep{ching2018coxnnet} &
Ten TCGA cohorts &
Gene expression &
Cox-nnet &
Concordance index \\

Hao et al. \citep{hao2019coxpasnet} &
TCGA brain and ovarian &
Expression, clinical, pathway &
Cox-PASNet &
Concordance index \\

Chaudhary et al. \citep{chaudhary2018deep} &
TCGA liver, $n = 360$ &
Expression, microRNA, methylation &
Autoencoder &
Concordance index 0.68; independent cohorts \\

Poirion et al. \citep{poirion2021deepprog} &
32 TCGA cancers &
Multiple data types &
DeepProg &
Concordance index; independent cohorts \\

Lee et al. \citep{lee2020dynamic} &
Cystic fibrosis registry, $n = 5{,}883$ &
Repeated clinical measurements &
Dynamic-DeepHit &
Time-dependent discrimination \\

\midrule
This work &
Four TCGA cancers, $n = 1{,}620$ &
Expression and clinical &
Cox elastic net &
Held-out test split; concordance index, area under the curve, Brier score \\
\bottomrule
\end{tabular}%
}
\end{table}

\subsection{Limitations of existing approaches}

Several limitations recur across this literature. Sample sizes are small
relative to the number of measurements, which raises the risk of unstable
weights. TCGA records one tissue sample per patient at diagnosis, so models that
need repeated measurements over time cannot be applied directly
\citep{lee2020dynamic}. Forcing a time-varying structure onto single-timepoint
data requires carrying values forward artificially, which introduces immortal
time bias \citep{suissa2008immortal}. Learned variables from autoencoders are
hard to map back to named genes \citep{chaudhary2018deep}. Pooling several
cancer types into one model averages gene effects that may differ between those
types. Most importantly, discrimination is reported far more often than
calibration or per-patient uncertainty, and external validation remains
inconsistent.

\subsection{Research gap}

Penalised Cox regression \citep{tibshirani1997lasso, simon2011regularization},
variable compression \citep{vanhouwelingen2006crossvalidated}, probability
calibration \citep{niculescu2005predicting}, uncertainty estimation
\citep{lee2020dynamic} and gene-level explanation \citep{hao2019coxpasnet} have
each been studied separately. No published pipeline combines all five at once. Such a pipeline must
also keep the training, calibration and test data strictly separated. It must
return an exact gene-level explanation rather than an approximate one. That
combination is the contribution of this work. The current model uses one measurement per patient,
taken at diagnosis. A version using repeated measurements is left for future
work.

% =============================================================================
\section{Proposed Method}
% =============================================================================

\subsection{Cohort assembly}

Patients were drawn from four TCGA cohorts. The first was glioblastoma
multiforme (TCGA study code GBM), a brain cancer \citep{brennan2013somatic}. The
second was liver hepatocellular carcinoma (TCGA study code LIHC)
\citep{tcga2017hepatocellular}. The third was pancreatic adenocarcinoma (TCGA
study code PAAD) \citep{tcga2017pancreatic}. The fourth was skin cutaneous
melanoma (TCGA study code SKCM) \citep{tcga2015melanoma}. The codes GBM, LIHC,
PAAD and SKCM are dataset
identifiers used by TCGA, not standard clinical abbreviations. Clinical records
and RNA-Seq gene activity values were downloaded in the file format published by
cBioPortal \citep{cerami2012cbio, gao2013integrative}. The four cohorts were
merged into one table of $N = 1{,}620$ patients.

Clinical fields that exist in only one cohort were removed. The melanoma skin
type field is one example. Such a field is missing for every patient in three of
the four cohorts. Filling those gaps with a single value would create an
artificial cluster of identical entries.

The cohort was split into three parts. Training used 972 patients, calibration
used 324 and testing used 324. The split was stratified on whether a death was
observed, rather than on survival time. Stratifying on survival time under right
censoring shifts the distribution between splits. Stratifying on the observed
event instead held the death rate nearly constant: 58.4\% in training, 58.6\% in
calibration and 58.3\% in testing.

\subsection{Removing redundant genes}

The 500 candidate gene activity columns were screened before any scaling or
compression. The absolute Pearson correlation was computed between every pair of
those 500 columns, using the training split only. Columns were then examined in
a fixed order. A column was discarded if its absolute correlation with any
earlier column exceeded 0.75. This rule keeps the first member of each
correlated group and discards the rest.

On this cohort, 36,205 gene pairs exceeded the threshold. The filter discarded
441 of the 500 genes and kept 59. Discarded genes were not arbitrary. The
strongest pairs were genes known to act together. The two acyl-coenzyme A
synthetase genes ACSM2A and ACSM2B correlated at $r = 0.996$. The three
fibrinogen chain genes FGA, FGB and FGG correlated above $r = 0.98$ with each
other. This pattern indicates that the filter removed genuine biological
redundancy rather than random noise. The screen costs one correlation matrix. It
avoids the repeated matrix inversions that a full variance inflation factor
audit would require. Figure~\ref{fig:gram_schmidt_collinearity} shows the
correlation structure of a representative subset of genes before and after this
step.

\begin{figure}[H]
    \centering
    \begin{subfigure}{0.48\linewidth}
        \centering
        \includegraphics[width=\linewidth]{images/fig3a_collinearity_before.png}
        \caption{Before filtering}
    \end{subfigure}
    \hfill
    \begin{subfigure}{0.48\linewidth}
        \centering
        \includegraphics[width=\linewidth]{images/fig3b_collinearity_after.png}
        \caption{After filtering}
    \end{subfigure}
    \caption{\small Pairwise correlation between a representative subset of
    genes, measured on the training split. Dark cells on the left mark groups of
    genes that move together, such as the fibrinogen chain genes. On the right,
    the discarded genes are marked and only weakly correlated genes remain.}
    \label{fig:gram_schmidt_collinearity}
\end{figure}

\subsection{Scaling and compression}

RNA-Seq values are counts. Their distribution has a long right tail, because a
few genes are far more active than the rest \citep{love2014moderated}. A
base-two logarithm, $\log_2(x + 1)$, was applied to the 59 surviving columns to
shorten that tail. The added one prevents an undefined logarithm where activity
is zero.

The transformed values were then scaled to have a mean of zero and a standard
deviation of one. This produced the matrix
$Z_{\text{gene}} \in \mathbb{R}^{N \times 59}$. The scaling was computed inside
each cross-validation training fold only. Computing it on all data first would
leak information from the validation fold into training.

Scaling has to happen before compression. PCA finds the directions of largest
variation \citep{hotelling1933analysis, jolliffe2016principal}. Without scaling,
a gene whose numbers simply span a wider range would dominate the first
direction, whatever its biological role. The loading matrix
$V \in \mathbb{R}^{59 \times 50}$, which satisfies $V^{T} V = I_{50}$, was
computed by singular value decomposition. The scaled genes were projected onto
$K = 50$ summary variables, giving
$X_{\text{pca}} = Z_{\text{gene}} V \in \mathbb{R}^{N \times 50}$.
Figure~\ref{fig:pca_gene_loadings} shows how the surviving genes map onto the
leading summary variables.

\subsection{Model input and objective}

Clinical fields were processed on a separate branch. Sex, race, ethnicity,
cancer type and age group were one-hot encoded, which turns each category into
its own binary column. Age was filled with the median where missing, then
scaled. This produced a clinical block
$X_{\text{clin}} \in \mathbb{R}^{N \times 20}$. The clinical block and the
50 summary variables were joined and rescaled, giving the model input
$X = [X_{\text{clin}} \mid X_{\text{pca}}] \in \mathbb{R}^{N \times 70}$.

The Cox model assigns each patient $i$ a risk score $\eta_i = X_i \beta$, with
hazard $h(t \mid X_i) = h_0(t) \exp(\eta_i)$ \citep{cox1972regression}. The
weights are $\beta = [\beta_{\text{clin}} \mid \beta_{\text{pca}}] \in
\mathbb{R}^{70}$. Fitting requires the Breslow partial log-likelihood
\citep{breslow1972discussion}. Evaluated directly, this needs a nested loop over
patients at every observed death, which costs $O(N^2)$ operations. Death times
were instead sorted in descending order, and running totals of the exponentiated
risk scores were accumulated in a single pass. This reduces the cost to $O(N)$.

An elastic-net penalty was added to the negative partial log-likelihood
\citep{zou2005regularization, wu2012elastic}. The penalty has an absolute-value
term and a squared term. The absolute-value term drives weights of uninformative
variables to zero. The squared term shrinks correlated variables together.

\subsection{Optimisation}

The objective was minimised with the limited-memory
Broyden-Fletcher-Goldfarb-Shanno algorithm with bound constraints (L-BFGS-B)
\citep{byrd1995limited}. This method needs a gradient everywhere. The
absolute-value term has no gradient at zero, and produced undefined values
during development. The absolute value was therefore replaced by the smooth
approximation $\sqrt{\beta^2 + \epsilon}$ with $\epsilon = 10^{-6}$. That value
is small enough to keep the sparsifying behaviour near zero. It is also far
enough above double-precision rounding error to keep the gradient stable.

The risk score $\eta_i$ was clipped to the range $[-40, 40]$ before
exponentiation. The upper bound gives $\exp(40) \approx 2.35 \times 10^{17}$,
which already represents an extreme relative risk. It stays well below the
double-precision overflow point at $\exp(709)$.

Two hyperparameters were tuned by grid search under five-fold nested
cross-validation. The penalty strength $\alpha$ was searched over
$\{0.1, 0.3, 0.8, 1.5, 3.0\}$. The balance between the two penalty terms,
$l1\_ratio$, was searched over $\{0.0, 0.1, 0.3, 0.5, 0.7\}$. The search
selected $\alpha = 3.0$ and $l1\_ratio = 0.7$. The final fit converged after 51
iterations. Figure~\ref{fig:inference_xai_architecture} shows how the fitted
weights and loadings feed the later stages. The implementation used scikit-learn
for preprocessing and cross-validation \citep{pedregosa2011scikit}.

\begin{figure}[H]
    \centering
    \includegraphics[width=0.9\linewidth]{images/fig2_inference_xai_architecture.png}
    \caption{Flow through the prediction and explanation stages. Fitted weights
    and compression loadings feed the risk score, the probability calibrator and
    the reversal step that produces gene-level and patient-level explanations.}
    \label{fig:inference_xai_architecture}
\end{figure}

\subsection{Converting risk scores into probabilities}

The risk score $\eta_i$ is only meaningful relative to other patients. It does
not answer the question a patient asks, which is the chance of surviving past a
given date. The 324-patient calibration split was held out before any filtering
or scaling took place. Isotonic regression was fitted on that split
\citep{ayer1955empirical, zadrozny2002transforming, niculescu2005predicting}.
Isotonic regression fits a step function that never decreases, so it cannot
reverse the ordering the Cox model already produced. Separate calibrators were
fitted at 12, 24, 36 and 60 months. Each maps a risk score to a survival
probability, using the observed survival fractions from the
Kaplan-Meier estimator as the target \citep{kaplan1958nonparametric}.
Figure~\ref{fig:isotonic_calibration_curves} plots the four fitted mappings.

\subsection{Estimating uncertainty by simulation}

A single predicted number gives no sense of the spread of possible outcomes.
Individual survival times were therefore simulated against the Breslow
cumulative baseline hazard $\Lambda_0(t)$, estimated from the training split
\citep{breslow1972discussion}. For each patient, 5,000 independent values
$U^{(k)} \sim \mathrm{Uniform}(0, 1)$ were drawn. Each draw was converted to a
target hazard, $\Lambda_{\text{target}} = -\ln(U^{(k)}) / \exp(\eta_i)$. The
matching month was located by binary search over the baseline hazard step
function. That lookup costs $O(\log E)$ operations, where $E = 461$ is the number
of distinct observed death times.

The 5,000 draws per patient were summarised as a pessimistic bound ($P10$), a
median ($P50$) and an optimistic bound ($P90$). The simulated survival curve was
also integrated up to 60 months to give the restricted mean survival time
\citep{royston2013restricted}. Integration was capped at the longest observed
follow-up in the cohort, 218.4 months, so the simulation never extrapolates
beyond the data. Figure~\ref{fig:monte_carlo_trajectories} shows example bands.

\subsection{Recovering gene-level weights exactly}

PCA reduced 59 genes to 50 summary variables, so the fitted weights attach to
those summary variables rather than to genes. Both the compression and the Cox
risk score are linear operations. The compression can therefore be reversed
exactly.

The sub-vector of weights attached to the summary variables is
$\beta_{\text{pca}} \in \mathbb{R}^{50}$. Multiplying the loading matrix by this
sub-vector gives a risk weight for every surviving gene:

\begin{equation}
W_{\text{gene}} = V \cdot \beta_{\text{pca}} \in \mathbb{R}^{59}.
\end{equation}

For patient $i$ and gene $g$, the contribution to the risk score follows
directly:

\begin{equation}
\Delta \eta_{g,i} = Z_{g,i} \cdot W_{\text{gene},g}.
\end{equation}

Here $Z_{g,i}$ is the scaled activity value that entered the compression step.
Because this is exact linear algebra, the contributions of all 70 model inputs
add up to the patient's risk score with no residual. Sorting them by size
produces a breakdown that can be checked term by term, shown in
Figure~\ref{fig:xai_plots}. Sampling-based methods such as SHAP
\citep{lundberg2017unified} and LIME \citep{ribeiro2016why} estimate the same
quantity, but only approximately.

\begin{figure}[H]
    \centering
    \includegraphics[width=\linewidth]{images/fig4_pca_gene_loadings.png}
    \caption{Loadings of a representative subset of the 59 surviving genes on
    the first ten summary variables. Cell values give the size of each gene's
    contribution to the corresponding summary variable.}
    \label{fig:pca_gene_loadings}
\end{figure}

% =============================================================================
\section{Results}
% =============================================================================

\subsection{Feature reduction}

The model input shrank in stages. The pipeline started with 506 candidate
columns, made up of 500 gene activity columns and 6 clinical fields. The
correlation filter left 65 columns, of which 59 were genes and 6 were clinical.
One-hot encoding then expanded the 6 clinical fields into 20 columns, while PCA
compressed the 59 genes into 50 summary variables. The final model input had 70
columns. Table~\ref{tab:dim_progression} lists these stages.

The correlation filter removed 441 of 500 genes, which is 88.2\% of the panel.
This indicates that most of an unfiltered 500-gene panel carries repeated
information rather than independent signal. It supports filtering before
compressing, instead of relying on PCA alone to absorb the redundancy.

\begin{table}[H]
\centering
\caption{Number of columns at each stage of the pipeline.}
\label{tab:dim_progression}
\begin{tabular}{lc}
\toprule
Stage & Columns \\
\midrule
Candidate gene activity columns & 500 \\
Clinical fields, before encoding & 6 \\
Gene columns after correlation filter & 59 \\
Clinical fields after correlation filter & 6 \\
Summary variables after compression & 50 \\
Clinical columns after one-hot encoding & 20 \\
Final model input $X$ & 70 \\
\bottomrule
\end{tabular}
\end{table}

\subsection{Model performance}

Five-fold nested cross-validation on the training split selected $\alpha = 3.0$
and $l1\_ratio = 0.7$. The mean cross-validated concordance index was 0.722,
with a standard deviation of 0.021 across folds. The concordance index measures
how often the model ranks two patients in the correct survival order
\citep{harrell1982evaluating}. A value of 0.5 matches random ordering, and 1.0
matches perfect ordering. Censoring-adjusted variants of this statistic have also
been proposed \citep{uno2011c}.

The final model reached a concordance index of 0.749 on the training split and
0.746 on the calibration split. On the held-out test split of 324 patients it
reached 0.736. Table~\ref{tab:cindex} lists these values. The gap between training
and test performance was 0.013. A gap this small indicates limited overfitting,
which is consistent with the aggressive filtering and shrinkage applied earlier.

\begin{table}[!htbp]
\centering
\caption{Concordance index by data split.}
\label{tab:cindex}
\begin{tabular}{lc}
\toprule
Split & Concordance index \\
\midrule
Training ($n = 972$) & 0.749 \\
Calibration ($n = 324$) & 0.746 \\
Test ($n = 324$) & 0.736 \\
Five-fold nested cross-validation (mean $\pm$ s.d.) & $0.722 \pm 0.021$ \\
\bottomrule
\end{tabular}
\end{table}

The fitted weights show which inputs mattered most. Cancer type carried the
largest clinical effect. Membership of the brain cancer cohort gave the largest
risk-increasing weight ($\beta = 0.772$). Membership of the melanoma cohort gave
the largest protective weight ($\beta = -0.322$). This ordering matches the
observed survival in these cohorts, reported in
Table~\ref{tab:cancer_outcomes}. Patient age gave the third-largest clinical
weight ($\beta = 0.293$), increasing risk. The first summary variable gave the
largest weight among the compressed gene variables ($\beta = 0.304$).

The two blocks contributed comparable amounts. Summed across all 50 summary
variables, the absolute weights of the gene block totalled 2.814. The 20
clinical columns totalled 2.086. Both blocks were scaled the same way, so these
totals are directly comparable. By this measure the compressed gene signal
carried slightly more weight than the full clinical panel.
Figure~\ref{fig:xai_plots}(a) ranks the ten largest weights across both blocks.

\subsection{Calibration at fixed time horizons}

The calibrators were evaluated on the 324-patient test split. At each horizon
the evaluation was restricted to patients whose outcome at that horizon is
known, after accounting for right censoring. That gave 289, 247, 228 and 218
patients at 12, 24, 36 and 60 months. Table~\ref{tab:horizon_metrics} reports
the area under the curve and the Brier score at each horizon.

The area under the curve measures discrimination at a fixed time point
\citep{heagerty2005survival}. It rose from 0.759 at 12 months to a peak of 0.860
at 24 months, then settled near 0.847 at 36 and 60 months. The Brier score
measures the squared difference between the predicted probability and the
observed outcome, so lower values are better \citep{graf1999assessment}. It fell
from 0.178 at 12 months to 0.132 at 60 months.

The 60-month horizon had both the lowest error and the highest observed death
rate, at 78.4\%. The death rate rises with horizon length for a mechanical
reason. As the window extends, fewer patients remain who were censored before
the horizon. More of the uncertainty at 12 months therefore comes from patients
last seen alive and event-free. Figure~\ref{fig:isotonic_calibration_curves}
plots the fitted mapping at each horizon. Each curve rises without ever falling,
which confirms that the calibrator preserved the ordering the Cox model
produced.

\begin{table}[!htbp]
\centering
\caption{Calibration results on the held-out test split.}
\label{tab:horizon_metrics}
\begin{tabular}{ccccc}
\toprule
Horizon (months) & Patients with known outcome & Death rate & Area under curve & Brier score \\
\midrule
12 & 289 & 0.284 & 0.759 & 0.178 \\
24 & 247 & 0.543 & 0.860 & 0.152 \\
36 & 228 & 0.662 & 0.847 & 0.150 \\
60 & 218 & 0.784 & 0.847 & 0.132 \\
\bottomrule
\end{tabular}
\end{table}

\begin{figure}[H]
    \centering
    \includegraphics[width=\linewidth]{images/fig5_isotonic_calibration_curves.png}
    \caption{Fitted calibration curves at (a) 12 months, (b) 24 months,
    (c) 36 months and (d) 60 months. Each panel compares the uncalibrated Cox
    risk score with the calibrated survival probability on the test split. The
    diagonal marks perfect calibration.}
    \label{fig:isotonic_calibration_curves}
\end{figure}

\subsection{Simulated uncertainty}

The simulation produced individual $P10$, $P50$ and $P90$ survival bounds for
all 324 test patients. The mean width of the interval between $P10$ and $P90$
was 107.8 months. The median width was 103.0 months. Intervals this wide reflect
the difficulty of projecting a single tissue sample forward across a follow-up
window reaching 218.4 months.

Interval width varied with risk. One high-risk test patient with
$\eta_i = +0.965$ received bounds of $P10 = 3.3$, $P50 = 13.6$ and
$P90 = 38.1$ months, with a 60-month restricted mean survival time of 16.9
months. One low-risk test patient with $\eta_i = -1.578$ received much wider
bounds of $P10 = 20.8$ and $P50 = 152.2$ months. That patient's $P90$ reached
the 218.4-month cap set by the longest observed follow-up, and the restricted
mean survival time was 52.3 months.

This asymmetry follows from the exponential form of the hazard. A large positive
risk score compresses the simulated death times towards the present, whatever
value $U^{(k)}$ takes. A large negative score spreads the same distribution
across the whole follow-up range.
Figure~\ref{fig:monte_carlo_trajectories} plots both cases. In practice the
$P10$ bound supplies a worst-case timeline for planning aggressive treatment.
The $P90$ bound supplies a best-case ceiling that is limited by the cohort's own
follow-up rather than by unbounded extrapolation.

\begin{figure}[H]
    \centering
    \includegraphics[width=\linewidth]{images/fig6_monte_carlo_trajectories.png}
    \caption{Simulated survival probability over a 60-month window for one
    high-risk and one low-risk test patient.}
    \label{fig:monte_carlo_trajectories}
\end{figure}

\subsection{Gene-level and patient-level explanations}

Reversing the compression resolved the first summary variable into its
constituent genes. The ten genes loading most strongly on that variable were
CD24, CXCL14, GPRC5A, PPP1R1B, SYT13, SERPINA3, CA9, ADAM6, FGFR2 and CXCL5.
Two of these have documented links to low-oxygen tumour regions and to spread
beyond the original site: CA9 \citep{giatromanolaki2001carbonic} and CXCL5
\citep{zhou2012overexpression}.

At the level of one patient, the breakdown splits the risk score across all 70
model inputs with no residual. Figure~\ref{fig:xai_plots}(b) shows a
representative test patient, identifier TCGA-DD-AACJ, who was still alive at
last contact after 69.0 months of follow-up. For this patient the largest single
term was protective: not belonging to the brain cancer cohort, contributing
$-0.616$. The two largest risk-increasing terms were patient age, contributing
$+0.328$, and one summary variable, contributing $+0.321$. All 70 terms summed
to the patient's risk score of $-0.028$. Because this decomposition is exact
matrix algebra, the numbers can be audited individually.

\begin{figure}[H]
    \centering
    \begin{subfigure}{0.48\linewidth}
        \centering
        \includegraphics[width=\linewidth]{images/fig7a_global_gene_importance.png}
        \caption{Largest fitted weights}
    \end{subfigure}
    \hfill
    \begin{subfigure}{0.48\linewidth}
        \centering
        \includegraphics[width=\linewidth]{images/fig7b_patient_risk_waterfall.png}
        \caption{Risk breakdown for one patient}
    \end{subfigure}
    \caption{Explanations produced by reversing the compression step.
    (a) Inputs ranked by fitted weight across the whole cohort.
    (b) Additive breakdown of one patient's risk score.}
    \label{fig:xai_plots}
\end{figure}

\subsection{Independent statistical screen}

The raw clinical and gene activity columns were also tested for association with
mortality, separately from the model. Two tests were used. A two-sided
Mann-Whitney U test compared each variable against whether a death was observed
\citep{mann1947test}. A Spearman rank correlation compared each variable against
survival time in months \citep{spearman1904proof}. Both were corrected for
multiple testing across all 1,049 evaluated columns of the 1,620-patient cohort,
using the Benjamini-Hochberg procedure \citep{benjamini1995controlling}. Both
tests examine one variable at a time, so they show association rather than an
effect independent of the other variables.

Cancer type was strongly associated with mortality
($p = 1.09 \times 10^{-53}$, corrected $q = 1.36 \times 10^{-51}$, rank-biserial
effect size 0.392). This matches the median survival across the pooled cohort,
which ranged from 12.16 months for brain cancer to 35.78 months for melanoma
(Table~\ref{tab:cancer_outcomes}). Tumour stage recorded under the American
Joint Committee on Cancer system was also strongly associated
($p = 8.74 \times 10^{-45}$, corrected $q = 5.15 \times 10^{-43}$). Stage
therefore retains predictive value even inside a pooled analysis of four cancer
types.

Among individual gene activity columns, DKFZp547D155 showed the strongest
correlation with survival time ($\rho = -0.346$, corrected
$q = 7.57 \times 10^{-31}$). Several genes recovered by reversing the
compression also passed this screen independently. CA9 gave $\rho = -0.338$
(corrected $q = 7.23 \times 10^{-30}$) and CXCL5 gave $\rho = -0.308$
(corrected $q = 4.33 \times 10^{-25}$). The screen makes no use of the fitted
model, so this agreement provides a check on the model-derived gene ranking.

\begin{table}[!htbp]
\centering
\caption{Outcome by cancer type across the pooled cohort ($N = 1{,}620$).}
\label{tab:cancer_outcomes}
\begin{tabular}{llccc}
\toprule
Cancer type & TCGA code & $n$ & Death rate & Median survival (months) \\
\midrule
Glioblastoma multiforme      & GBM  & 595 & 0.827 & 12.16 \\
Pancreatic adenocarcinoma    & PAAD & 184 & 0.543 & 15.34 \\
Liver hepatocellular carcinoma & LIHC & 372 & 0.355 & 19.76 \\
Skin cutaneous melanoma      & SKCM & 469 & 0.475 & 35.78 \\
\bottomrule
\end{tabular}
\end{table}

% =============================================================================
\section{Discussion}
% =============================================================================

\subsection{Why a linear model was chosen}

A neural network could process all 500 candidate genes without a correlation
filter or a compression step. On a cohort of this size that flexibility is
expensive. With 972 training patients, a network large enough to capture
interactions between genes has more trainable parameters than it has training
examples. The usual defence is heavy regularisation, which in practice can push
the network back towards a linear function.

The model reported here keeps a linear structure from the start. The correlation
filter removed 441 redundant genes before they could destabilise the fitted
weights. The smoothed elastic-net penalty then shrank the remaining correlated
variables together, instead of selecting one representative from each group
\citep{zou2005regularization}. The resulting test concordance index of 0.736 is
comparable to the 0.68 reported by an autoencoder model on 360 TCGA liver cancer
patients \citep{chaudhary2018deep}. The two numbers come from different cohorts
and different splits, so this is a reference point rather than a controlled
comparison.

Reversing the compression extends readability past the summary variables and
back onto named genes. That computation is exact. Sampling-based explainers
estimate the same quantity approximately \citep{lundberg2017unified,
ribeiro2016why}. An exact decomposition is useful wherever the same
explanation must be reproducible across repeated runs.

\subsection{Clinical readability}

A relative risk score can tell a clinician that one patient is at higher risk
than another. It cannot answer what that patient's survival probability is over
a stated period. The calibration layer supplies that conversion, with Brier
scores between 0.132 and 0.178 across the four horizons. The simulation layer
then converts a probability into a range of plausible survival months.

The asymmetry of those ranges carries information. Bounds were tight for
high-risk patients and wide for low-risk patients. This pattern suggests the
model is confident about near-term decline in aggressive cases, and
appropriately uncertain about long-horizon outcomes in more favourable cases.

Simulation was used rather than a confidence interval on the fitted weights. A
confidence interval on $\beta$ describes a population-level parameter. The
simulation estimates a distribution for one patient's survival time. Because the
baseline hazard step function ends at the cohort's longest observed follow-up of
218.4 months, the simulation cannot project beyond the available data.

\subsection{Limitations}

Six limitations should be stated.

The proportional hazards assumption was not tested. Schoenfeld residuals provide
the standard check \citep{schoenfeld1982partial}, and were not computed. The
assumption is likely violated for at least some of the four cancer types, given
their different progression timelines.

Both selected hyperparameters sat at the upper edge of their search grids.
$\alpha = 3.0$ was the largest penalty strength searched, and
$l1\_ratio = 0.7$ was the largest balance value searched. A stronger penalty or
a more lasso-weighted balance may perform better and was not explored.

Pooling four histologically different cancers into one weight vector averages
gene effects across them. A gene that is protective in one cancer type and
harmful in another will have those effects cancel. Subtype-specific precision is
traded for coverage across cancer types.

The model has not been tested outside TCGA. Any use beyond this retrospective
setting would need testing against batch effects and different sequencing
platforms \citep{poirion2021deepprog}.

Each patient is represented by one tissue sample taken at diagnosis. Disease
progression over time is therefore not modelled. Methods that use repeated
measurements need data TCGA does not provide \citep{lee2020dynamic}.

The statistical screen tested one variable at a time. It shows association with
mortality and does not establish that any variable adds predictive value beyond
the others.

% =============================================================================
\section{Conclusion}
% =============================================================================

A Cox proportional hazards model with an elastic-net penalty was fitted to
clinical records and gene activity measurements from four pooled TCGA cancer
cohorts. Redundant genes were removed first, reducing 500 candidates to 59, and
those were compressed into 50 summary variables. The resulting 70-column model
reached a concordance index of 0.736 on a held-out test split of 324 patients.
Isotonic regression converted the raw risk scores into survival probabilities,
with Brier scores as low as 0.132. A 5,000-draw simulation against the Breslow
baseline hazard turned those probabilities into per-patient ranges of plausible
survival months. Reversing the compression by matrix multiplication returned an
exact risk weight for each of the 59 genes, so no approximate explainer was
needed. A separate single-variable screen across all 1,049 columns recovered
several of the same genes, including CA9 and CXCL5. Three steps are needed next.
Interaction terms specific to each cancer type should be added. The pipeline
should be tested on cohorts outside TCGA. The proportional hazards assumption
should be checked with Schoenfeld residuals before any clinical use is
considered.

% =============================================================================
\section*{Acknowledgement}

[Insert acknowledgement here.]

% ---------------------------------------------------------------- refs -------
\vspace{1em}
\bibliographystyle{unsrtnat}
\bibliography{references}

\end{document}
